%% =====================================================================
\section{Literature Review}
\label{sec:related}
%% =====================================================================

We organize prior work into six threads: cache stampedes and large-scale caching;
denial of service and application-layer attacks; attacks that target caching
infrastructure; availability attacks in cloud and container environments; overload
control; and detection. We close with a literature matrix and the research gap.

\subsection{Cache Stampede and the Thundering Herd}
Fitzpatrick's description of memcached established the cache-aside pattern that most
web systems still follow \cite{fitzpatrick2004}. Nishtala et al., describing
Facebook's memcache deployment, identify ``stale sets'' and ``thundering herds'' as
first-order concerns and introduce \emph{leases}: a per-key token that permits only
one client to regenerate a missing value while others briefly wait or accept slightly
stale data. They report that leases substantially reduce peak database query rates
\cite{nishtala2013}. Vattani, Chierichetti, and Lowenstein formalize the mitigation
side and derive an optimal \emph{probabilistic early recomputation} rule (XFetch) in
which an item is refreshed before its TTL with a probability that rises as expiry
approaches \cite{vattani2015}. Workload studies of large key--value stores
\cite{atikoglu2012} and of photo caching at scale \cite{huang2013} confirm the
heavy skew that makes hot-key protection critical, and Fan et al.\ show that a small
cache in front of a partitioned service provably balances load only if it holds the
hottest keys \cite{fan2011}.

All of these works treat the stampede as a \emph{performance} hazard arising from
benign load. CacheCollapse builds directly on their mitigation insights---single-flight
regeneration is the descendant of memcache leases---but reframes the stampede as an
\emph{attacker-triggerable} availability vulnerability, adds detection, and measures
end-to-end user-visible availability rather than database load alone.

\subsection{Denial of Service and Application-Layer Attacks}
Classical DDoS research classifies attacks and defenses mainly at the network and
transport layers. Mirkovic and Reiher give the canonical taxonomy of attacks and
defenses \cite{mirkovic2004}, and Zargar et al.\ survey flooding defenses, stressing
that protection must span prevention, detection, and response \cite{zargar2013}.
Application-layer DoS is distinct and harder. Xie and Yu model normal browsing behavior
to detect attacks composed of ordinary HTTP GET requests \cite{xie2009}; Ranjan et al.\
propose DDoS-Shield, which assigns suspicion scores to sessions and schedules requests
to limit the damage of attackers mimicking legitimate clients \cite{ranjan2009};
and Praseed and Thilagam survey the space, emphasizing that such attacks can be
mounted at low volume with legitimate requests \cite{praseed2019}.

A central difficulty is distinguishing attacks from \emph{flash crowds}---sudden,
legitimate surges of interest. Jung et al.\ characterize flash crowds and DoS attacks
and propose distinguishing features such as client distribution and per-client request
rates \cite{jung2002}; Kandula et al.\ address attacks that deliberately mimic flash
crowds by using CAPTCHAs and admission control to separate humans from bots
\cite{kandula2005}. Two further lines show that damaging attacks need not be
voluminous: low-rate ``shrew'' attacks exploit TCP's retransmission timer with
carefully timed bursts \cite{kuzmanovic2003}, and algorithmic-complexity attacks feed
inputs that trigger the worst-case behavior of legitimate data structures
\cite{crosby2003}. The cache-stampede IDDoS shares the timing element of the former and
the ``legitimate but expensive path'' element of the latter: requests are legitimate,
the aggregate web-tier rate need not be anomalous, and the damage is produced
downstream by cache semantics.

\subsection{Attacks that Target Caching Infrastructure}
A growing body of work treats caches themselves as attack surfaces. Triukose et al.\
show that content-delivery networks can be abused to amplify load on an origin server
by requesting uncacheable or cache-busting URLs \cite{triukose2009}. Guo et al.\
demonstrate that CDN forwarding behavior can be turned against the very origin it is
meant to protect (``CDN Judo'') \cite{guo2020}. Nguyen et al.\ introduce
\emph{cache-poisoned denial of service} (CPDoS), in which an attacker causes an
intermediate HTTP cache to store an error page that is then served to every subsequent
user \cite{nguyen2019}. These works share our view that caching changes the threat
model, but they target shared HTTP caches at the network edge and the integrity of
cached content. CacheCollapse targets the \emph{application-internal} cache and the
database behind it, and its attack requires neither cache-busting parameters nor
poisoned content---only correct timing.

\subsection{Availability Attacks in Cloud and Containerized Environments}
Containers are the dominant packaging and orchestration unit for modern applications
\cite{pahl2015,merkel2014,burns2016}, and the microservice style decomposes
applications into many interacting services \cite{dragoni2017}. Security analyses of
Docker catalog isolation and configuration weaknesses \cite{combe2016,sultan2019}, but
availability has received comparatively less attention. The most directly relevant
line is the \emph{Yo-Yo attack}, in which periodic bursts drive an auto-scaling
controller to oscillate between scale-up and scale-down, degrading performance and
inflating cost while evading simple volume detection \cite{sides2016,bremlerbarr2017}.
CacheCollapse is complementary: like Yo-Yo it exploits an emergent property of a
cloud-native pattern (there, elasticity; here, caching) using legitimate-looking
traffic, and like Yo-Yo it must be studied in an orchestrated, resettable environment.
Unlike Yo-Yo, it causes damage without any auto-scaler and in a single deployment.

\subsection{Overload Control and Graceful Degradation}
Systems research has long argued that services should degrade gracefully under
overload. SEDA structures a server as stages connected by queues with admission
controllers so that overload is shed explicitly \cite{welsh2001}, and Welsh and Culler
later propose adaptive overload control driven by measured response times
\cite{welsh2003}. Session-based admission control rejects new sessions under load to
preserve service for sessions already in progress \cite{cherkasova2002}. Practitioner
literature codifies related stability patterns---timeouts, circuit breakers, bulkheads,
and serving stale data---\cite{nygard2018}, and site-reliability practice treats load
shedding, client throttling, and protection against cache loss as defenses against
cascading failure \cite{beyer2016}. Rate limiting may be enforced at a single proxy or
globally across replicas \cite{raghavan2007}. A contribution of the present paper is to
show empirically that these controls are \emph{not} interchangeable in the stampede
setting: coalescing protects the database, edge rate limiting protects the request
path, and only their combination preserves end-to-end availability.

\subsection{Detection: Rules, Statistics, and Machine Learning}
Denning's intrusion-detection model introduced statistical profiles of normal behavior
as a basis for detecting attacks \cite{denning1987}. Chandola et al.\ and Hodge and
Austin survey anomaly detection broadly, underlining that the choice of method
depends on label availability, data type, and interpretability requirements
\cite{chandola2009,hodge2004}. Two cautionary results shape our evaluation. Axelsson
shows that because attacks are rare, even a small false-positive rate can make most
alarms false \cite{axelsson2000}; and Sommer and Paxson argue that machine-learning
detectors face a ``closed world'' problem in which anything not seen in training is
flagged, so that operational benign changes look like attacks \cite{sommer2010}. Among
unsupervised methods the Isolation Forest \cite{liu2008,liu2012} is widely used because
it needs no labels and few samples. We therefore pair it with an explainable rule and
evaluate both not only on attack and normal traffic but also on idle and defended
traffic, where the closed-world problem becomes visible. ROC analysis \cite{fawcett2006}
is used to separate the ranking quality of a score from the choice of threshold.

\subsection{Literature Matrix and Research Gap}
Table~\ref{tab:litmatrix} summarizes the most relevant works, what they contribute, and
what they leave open relative to the problem studied here.

\begin{table}[p]
\centering
\caption{Literature matrix: focus, approach, and gap relative to CacheCollapse.}
\label{tab:litmatrix}
\footnotesize
\begin{tabularx}{\linewidth}{@{}L{2.7cm}L{2.6cm}YY@{}}
\toprule
Work & Focus & Approach / key finding & Gap relative to this paper \\
\midrule
Nishtala et al.\ 2013 \cite{nishtala2013} & Caching at scale & Leases coalesce regeneration; stale sets & Performance framing; no adversary, no detection \\
Vattani et al.\ 2015 \cite{vattani2015} & Stampede prevention & Optimal probabilistic early recomputation (XFetch) & Assumes benign load; does not handle a burst that is already in flight \\
Fan et al.\ 2011 \cite{fan2011} & Load balancing & Small front cache balances partitioned back ends & Capacity planning, not attack/defense \\
Mirkovic \& Reiher 2004 \cite{mirkovic2004} & DDoS taxonomy & Classifies attacks and defenses & Network-layer focus \\
Zargar et al.\ 2013 \cite{zargar2013} & DDoS defense survey & Prevention--detection--response framing & Flooding attacks; no cache semantics \\
Xie \& Yu 2009 \cite{xie2009} & App-layer DDoS detection & Models normal browsing (HSMM) & Detection only; no mitigation; no cache \\
Ranjan et al.\ 2009 \cite{ranjan2009} & App-layer DDoS mitigation & Suspicion-based request scheduling & Assumes attacker sessions differ in behavior \\
Jung et al.\ 2002 \cite{jung2002}; Kandula et al.\ 2005 \cite{kandula2005} & Flash crowd vs.\ attack & Client distribution; CAPTCHA + admission control & Web-tier volume focus; not downstream amplification \\
Kuzmanovic \& Knightly 2003 \cite{kuzmanovic2003} & Low-rate DoS & Timed bursts exploit TCP timers & Transport layer \\
Crosby \& Wallach 2003 \cite{crosby2003} & Algorithmic complexity & Inputs trigger worst-case paths & Single-process data structures \\
Nguyen et al.\ 2019 \cite{nguyen2019}; Guo et al.\ 2020 \cite{guo2020} & Attacks on HTTP caches/CDNs & Cache poisoning; CDN abuse against origin & Shared edge caches; integrity or forwarding, not expiry timing \\
Sides et al.\ 2016 \cite{sides2016}; Bremler-Barr et al.\ 2017 \cite{bremlerbarr2017} & Cloud availability (Yo-Yo) & Bursts make auto-scaler oscillate & Targets elasticity, not caching; no layered mitigation \\
Combe et al.\ 2016 \cite{combe2016}; Sultan et al.\ 2019 \cite{sultan2019} & Container security & Isolation and configuration risks & No availability attack instrumented end to end \\
Welsh et al.\ 2001/2003 \cite{welsh2001,welsh2003} & Overload control & Staged queues; adaptive admission & General overload; no attack model \\
Liu et al.\ 2008/2012 \cite{liu2008,liu2012} & Anomaly detection & Isolation Forest & Method only; no domain evaluation \\
Sommer \& Paxson 2010 \cite{sommer2010} & ML for intrusion detection & Closed-world critique & Position paper; motivates our idle/defended evaluation \\
\midrule
\textbf{CacheCollapse} & Cache-stampede IDDoS & Containerized testbed + analytical model + rule \& ML detection + layered mitigation + ablation & Single-host, synthetic traffic (Sec.~\ref{sec:validity}) \\
\bottomrule
\end{tabularx}
\end{table}

\paragraph{Research gap.} Prior work offers (a)~mitigations for stampedes framed as a
performance problem, (b)~detectors for application-layer floods framed at the web
tier, and (c)~analyses of cache- and cloud-specific availability attacks at the network
edge or the auto-scaler. To our knowledge no prior study combines, in one reproducible
containerized environment, an \emph{adversarial} treatment of the application-internal
cache stampede, \emph{both} explainable and learned detection evaluated beyond the
attack/normal dichotomy, and an \emph{ablation} of layered mitigations measured at the
level of user-visible availability. CacheCollapse addresses that gap.
